% Recovered verbatim from chapters_16_20_with_part/chapters_16_20_with_part.tex
% Canonical chapter 19
\chapter{Complexity, Circuits, and Interior Growth}

\section{ER Networks as Quantum Circuits}

The ER network can be viewed as a graphical representation of a quantum circuit preparing the state of the universe from a simple reference.
Nodes correspond to gates or tensors, and ER links carry entangled degrees of freedom.
Circuit depth and gate counts define a notion of computational complexity for the state.

\section{Complexity and White-State Interiors}

For collapsing configurations, the growth of the interior region associated with a white-state core can be correlated with the growth of circuit complexity.
As the system evolves toward $\Wcore$, the complexity of the state (as measured along an appropriate circuit) increases, even as the classical geometry appears to ``collapse.''
This suggests a dual description in which interior volume or phase-space volume is related to computational complexity.

\section{Complexity Geometry}

We can formalise this by introducing a complexity metric on the space of unitaries acting on the combined Hilbert space, assigning costs to different generators corresponding to phase, chromatic, ER, and geometric transformations.
Geodesics in this metric approximate optimal circuits, and their lengths define state complexity.
In MK--43 Ultra, this complexity geometry is layered on top of the spacetime geometry encoded in $g_{\mu\nu}^{\mathrm{(eff)}}$.

